\chapter{Installation}
\label{ch:install}

Nothing in this chapter needs a robot. The result is a workspace that runs the
full simulation path.

\section{Prerequisites}

Ubuntu 24.04 Noble with ROS~2 Jazzy. The tested profile is exactly that
combination; nothing else has passed a gate.

\begin{lstlisting}
sudo apt install -y \
  ros-jazzy-desktop ros-jazzy-ros-gz ros-jazzy-gz-ros2-control \
  ros-jazzy-ros2-control ros-jazzy-ros2-controllers \
  ros-jazzy-moveit ros-jazzy-robot-state-publisher ros-jazzy-xacro \
  ros-jazzy-joint-state-publisher-gui \
  python3-colcon-common-extensions python3-rosdep python3-vcstool
\end{lstlisting}

\section{Source and the dependency pin}

One external repository is imported, at a full commit SHA:

\begin{lstlisting}
git clone https://github.com/yunusdanabas/baxter_ros2_jazzy.git
cd baxter_ros2_jazzy
source /opt/ros/jazzy/setup.bash
mkdir -p src
vcs import src < repos/baxter_core.repos
git -C src/baxter_common_ros2 rev-parse HEAD
# expect: 678bfabea8c895b4134951a6c076217a90b9e0e6
\end{lstlisting}

That import brings the CentraleNantes \texttt{baxter\_common\_ros2} message and
description packages. The pin is a full SHA rather than a branch on purpose: the
robot description and message definitions determine wire compatibility with a
physical robot, and a branch that moves underneath a hardware session is a defect
waiting for the worst possible moment.

If \texttt{src/baxter\_common\_ros2} already exists, verify the SHA rather than
re-importing over local changes.

\section{Build}

\begin{lstlisting}
rosdep install --from-paths src --ignore-src -r -y
colcon build --base-paths src --symlink-install --packages-skip baxter_bridge
source install/setup.bash
\end{lstlisting}

\texttt{--packages-skip baxter\_bridge} is not optional. That package links ROS~1
libraries which do not exist on a clean Jazzy machine; it belongs on a bridge host
with both distributions installed, and this workspace does not use it --- the
hardware path uses a pure-Python bridge instead (Chapter~\ref{ch:bridge}).

If the workspace or its underlay changes, remove \texttt{build/}, \texttt{install/}
and \texttt{log/} before rebuilding, and never hand-edit generated setup files.

\begin{safetybox}[The conda trap, which does not announce itself]
If a conda or mamba environment is active, its \texttt{python3} comes first on
\texttt{PATH}. \texttt{setuptools} bakes the active interpreter into the shebang
of every installed console script at \emph{build} time, so the build succeeds and
then every \texttt{ros2 run} fails at runtime with
\texttt{No module named 'rclpy.\_rclpy\_pybind11'}, because \texttt{rclpy}'s C
extension is built against the system interpreter.

Check before building, not after:

\begin{lstlisting}
conda deactivate
which python3      # want /usr/bin/python3
\end{lstlisting}

A successful build is not evidence that this is fine.
\end{safetybox}

\section{Verify}

\begin{lstlisting}
bash scripts/sim_smoke.sh gates
\end{lstlisting}

This runs the hardware-free checks: compile, lint, unit tests, the 25-case
hardware-bridge dry run, URDF expansion, and the MoveIt collision-matrix check.
It needs no Gazebo and no GPU. Chapter~\ref{ch:verify} describes what each check
is actually protecting.
